\subsection{Soft endpoints of the real-energy integrals}
The remaining rational and single-logarithm sectors have elementary
soft residues. Put $H=s-1$, and use $e=x$ at the first soft endpoint
and $e=y$ at the second. Here $e$ is an energy invariant, not the
charge. With $z=er$, the exact phase-space boundary is
\begin{equation}
 0<e<H,\qquad 0<r<\rho(e)=\frac{H-e}{H+1-e},\qquad
 xy=e[H+e(r-1)],\qquad \dd x\,\dd z=e\,\dd e\,\dd r.
 \label{eq:comp-real-soft-domain}
\end{equation}
To obtain this domain from $0<z<x(H-x)/(s-x)$, set
$x=e$ and divide the inequality by $e>0$. At the other soft
endpoint set $y=e$ and use $x=H-e+z$ in the same inequality.
Both give $r<\rho(e)$. The two changes of variables have Jacobian
magnitude $e$. A disjoint partition near the two soft corners
allows their residues to be added without double counting. Away
from both corners the regulators converge to one on compact
subsets, so their difference has zero limit there.

Let $q=\rho(0)=H/(H+1)$. For a sector with a simple soft
singularity, define its two residues by
\begin{equation}
 f_X(r)=\lim_{e\to0}e^2R(H,e,er),\qquad
 f_Y(r)=\lim_{e\to0}e^2R(H,H+e(r-1),er).
\end{equation}
The rational sector and the two terms multiplying $\ln s$ have
this behavior. The coefficient of $L_W$ is only of order $e^{-1}$,
and $L_W(0)$ is finite, so this sector needs no soft matching term.
All coefficients in this section precede the common normalization
$1/(8sH)$ of the real cross section.

Write $f=f_X+f_Y$. The regulator $(xy)^\eta$ gives a pole
$A/\eta$ and finite part $F_\eta$. If the energy primitives are
instead evaluated with $z>\zeta$, define their finite part by
subtracting $-A\ln\zeta$. The symbol $\operatorname{FP}$ denotes
the constant term after that subtraction. Then
\begin{align}
 A&=\int_0^q f(r)\,\dd r,\\
 F_\eta&=\operatorname{FP}_{\zeta\to0}
                \int_{z>\zeta}R\,\dd x\,\dd z
       +A\ln H-\int_0^q f(r)\ln r\,\dd r.
 \label{eq:comp-real-simple-matching}
\end{align}
Indeed, the soft measure is $f(r)\dd e/e$, while $z>\zeta$
sets its lower endpoint to $e=\zeta/r$. The analytic regulator
instead multiplies this measure by $(He)^\eta$ at leading order.
The difference produces the two logarithmic terms above. The
$O(e)$ change in $\rho(e)$ is integrable for a simple soft pole.

These residue integrals do not require numerical quadrature.
Partial fractions give a polynomial and poles at $r=1$ and
$r=-4/H$, both outside $0\leq r\leq q$. The necessary moments are
\begin{align}
 \int_0^q r^n\,\dd r&=\frac{q^{n+1}}{n+1},\\
 \int_0^q r^n\ln r\,\dd r
 &=q^{n+1}\left[\frac{\ln q}{n+1}-\frac1{(n+1)^2}\right],\\
 \int_0^q\frac{\dd r}{r-a}&=\ln(1-q/a),\\
 \int_0^q\frac{\ln r}{r-a}\,\dd r
 &=\ln q\ln(1-q/a)+\operatorname{Li}_2(q/a).
\end{align}
Here $n\geq0$ and $a$ is a pole outside the interval.
Repeated poles follow by applying
$\partial_a^{j-1}/(j-1)!$. Every logarithm above has a positive
argument, and $q/a<1$. These formulas therefore specify real
branches directly from the physical integration interval.
The independently generated residue coefficients and their exact
moments are recorded in
\nolinkurl{compton_real_soft_matching_integrated.json}.


\paragraph{Bose symmetry of the subtractions.} The two mixed angular
kernels $J_{UW}$ and $J_{VW}$ are exchanged by $x\leftrightarrow y$.
Their local Taylor subtractions, however, were expanded at fixed
$x$. Thus $T_{VW}(H,x,z)$ is not exactly
$T_{UW}(H,H-x+z,z)$. To replace the subtracted $VW$ sector by a
second copy of $UW$, the rational sector must receive
\begin{equation}
 D(H,x,z)=R_{VW}(H,x,z)
 \big[T_{UW}(H,H-x+z,z)-T_{VW}(H,x,z)\big].
 \label{eq:comp-real-bose-restoration}
\end{equation}
The bracket starts at $z^2$, so it introduces no artificial Gram
pole. It can nevertheless change the soft residue because its
expansion at fixed $x$ is not uniform as $y\to0$.

At this endpoint write
$D(H,H+e(r-1),er)=a(r)e^{-3}+b(r)e^{-2}+O(e^{-1})$.
The coefficients needed in this endpoint expansion are explicit:
\begin{align}
 a(r)&=\frac{16(H+1)(Hr+2)(Hr-2H-6)}{3H(r-1)^3},\\
 b(r)&=\frac{2(H+1)}{3H^2(r-1)^3(Hr+4)^2}
                    \sum_{j=0}^5 B_j(H)r^j.
\end{align}
The six polynomials are
\begin{align}
 B_0(H)&=256 H,\\
 B_1(H)&=64 \left(H^{3} + 6 H^{2} - 11 H - 12\right),\\
 B_2(H)&=16 \left(H^{4} + 2 H^{3} + 2 H^{2} + 128 H + 108\right),\\
 B_3(H)&=- 8 \left(H^{5} + 8 H^{4} - 19 H^{3} - 160 H^{2} - 192 H - 72\right),\\
 B_4(H)&=- 2 H \left(H^{5} + 3 H^{4} - 36 H^{3} - 108 H^{2} - 152 H - 96\right),\\
 B_5(H)&=H^{2} \left(H + 6\right) \left(H^{3} - 4 H^{2} - 12 H - 8\right).
\end{align}
These follow by substituting the exact soft coordinates into
\eqref{eq:comp-real-bose-restoration} and expanding the resulting
rational numerator and denominator through the first two powers.

The moving angular boundary now contributes to the logarithmic
soft coefficient. Since $\rho'(0)=-1/(H+1)^2$, it is
\begin{equation}
 A_D=\int_0^q b(r)\,\dd r-\frac{a(q)}{(H+1)^2}.
 \label{eq:comp-real-bose-pole}
\end{equation}
More explicitly, the radial density after the angular integration is
\begin{align}
 \int_0^{\rho(e)} eD\,\dd r
 &=\frac{1}{e^2}\int_0^q a(r)\,\dd r
   +\frac1e\left[\int_0^q b(r)\,\dd r+\rho'(0)a(q)\right]
   +O(1).
\end{align}
The first term is a power singularity, whose analytic continuation
has a pole at $\eta=1$. The coefficient of $1/e$ gives the pole
at $\eta=0$. This explains why expanding at a fixed angular upper
limit would miss the second term in
\eqref{eq:comp-real-bose-pole}. The terms generated by the
$O(\eta e)$ part of $(xy)^\eta$ affect its finite part and must
not be discarded.

Its finite contribution must also be included when combining the
real sectors. Omitting $D$ would violate the integrated soft
coefficient even though the original amplitude is Bose symmetric.

\paragraph{Integrated infrared check.} Denote the poles of the
rational, $\ln s$, $UV$, and $UW$ sectors by
$A_{\mathrm{rat}},A_{\log},A_{UV},A_{UW}$, respectively.
The required sum, before the common factor $1/(8sH)$, is
\begin{equation}
 A_{\mathrm{rat}}+A_{\log}+A_{UV}+2A_{UW}+A_D
 =8H\int_{-1}^1\mathcal B(c)\,[b_e(c)-1]\,\dd c.
 \label{eq:comp-real-integrated-pole-check}
\end{equation}
The right-hand side follows independently from the soft current
and the Born trace. It equals the coefficient required to cancel
the photon-mass logarithm of the virtual amplitude after the
regulator matching.
